% Research-Focused Academic CV Template
% Bundled with the academic-cv-editor skill.
% Personal data from the source template has been removed.
% Compile with pdfLaTeX / latexmk.

\documentclass[10pt,a4paper]{article}
\usepackage[T1]{fontenc}
\usepackage[utf8]{inputenc}
\usepackage{libertinus}
\usepackage[left=1.15cm,right=1.15cm,top=1cm,bottom=1cm]{geometry}
\usepackage{xcolor}
\usepackage{enumitem}
\usepackage{titlesec}
\usepackage{tabularx}
\usepackage{array}
\usepackage{setspace}
\usepackage[hidelinks]{hyperref}
\usepackage{microtype}

\pagestyle{empty}
\setlength{\parindent}{0pt}
\setlength{\parskip}{0pt}
\setstretch{1.05}
\urlstyle{same}

\definecolor{accent}{HTML}{1F355A}
\definecolor{textgray}{HTML}{20242A}
\color{textgray}

\titleformat{\section}
  {\fontsize{11pt}{11.8pt}\selectfont\bfseries\scshape\color{accent}}
  {}
  {0pt}
  {}
  [\vspace{-0.28em}\color{accent}\titlerule]
\titlespacing*{\section}{0pt}{7pt}{2.5pt}

\setlist[itemize]{
  leftmargin=1.35em,
  itemsep=1pt,
  topsep=1.7pt,
  parsep=0pt,
  partopsep=0pt
}

\newcommand{\datedline}[2]{%
  \begin{tabularx}{\textwidth}{@{}Xr@{}}
    #1 & \textit{#2}
  \end{tabularx}%
}

\newcommand{\roleentry}[3]{%
  \datedline{{\fontsize{10.3pt}{11.2pt}\selectfont\textbf{#1}}}{#3}\vspace{1.25pt}
  \textit{#2}\par\vspace{1.25pt}
}

\newcommand{\projectline}[1]{%
  \textbf{Project:} \textit{#1}\par\vspace{1.5pt}
}

\newcommand{\puzzleline}[1]{%
  \textbf{Research Puzzle:} #1\par\vspace{1.8pt}
}

\newcommand{\entryspace}{\vspace{3.75pt}}

\newcommand{\paperentry}[1]{%
  \hangindent=1.5em\hangafter=1 #1\par
}

% ---------- Identity: replace these values ----------
\newcommand{\cvname}{First Last}
\newcommand{\cvphone}{+00 000 0000 0000}
\newcommand{\cvemail}{name@example.edu}
\newcommand{\cvwebsite}{www.example.com}
\newcommand{\cvwebsiteurl}{https://www.example.com/}
\newcommand{\cvlocation}{City, Country}

\begin{document}
\fontsize{10.2pt}{11.3pt}\selectfont

\begin{center}
  {\fontsize{19pt}{20pt}\selectfont\bfseries \cvname}\par
  \vspace{3pt}
  {\fontsize{9.7pt}{10.4pt}\selectfont
  \cvphone \textbar\
  \href{mailto:\cvemail}{\cvemail} \textbar\
  \href{\cvwebsiteurl}{\cvwebsite} \textbar\ \cvlocation}
\end{center}
\vspace{-2pt}

% Remove any section not supported by the applicant's record.

\section{Research Profile}
Early-career researcher in [FIELD] with research interests in [SUBFIELD / PROBLEM]. Current research examines [RESEARCH FOCUS] using [METHODS / APPROACH].

\section{Education}
\datedline{\textbf{University Name} \textbar\ City, Country}{Expected Month YYYY}
\textit{Degree in Subject} \hfill \textbf{GPA / Classification: X}
\par\vspace{1.25pt}
\datedline{\textbf{Exchange / Visiting Institution} \textbar\ City, Country}{Term YYYY}
\textit{Programme / Visiting Status}\par

\section{Research Interests}
Subfield \textperiodcentered\ Research Problem \textperiodcentered\ Mechanism / Theme \textperiodcentered\ Regional or Empirical Focus

\section{Research Experience}
\roleentry
{Research Role}
{Institution \textbar\ Supervisor: Dr. Example Name}
{Mon--Mon YYYY}
\projectline{Research Project Title}
\begin{itemize}
    \item Analysed [evidence/material] to examine [research problem].
    \item Applied [method/approach] to [specific analytical task].
    \item Contributed to [conceptual, empirical, methodological, or presentation output].
\end{itemize}

\section{Selected Research Papers \& Working Papers}
\paperentry{
\textbf{Last, F.} ``Paper Title.''
\hfill \textit{Research Paper / Working Paper / Conference Paper, YYYY}\\[-1pt]
\textit{Keyword; keyword; keyword}\\[1.5pt]
\small
One or two lines explaining the research question, argument, cases, method, or evidence.\\[1.5pt]
\textbf{Status / Recognition:} Add only if factually supported.
}

\section{Publications}
\paperentry{Last, F. (YYYY). ``Publication Title.'' \textit{Venue}, volume(issue), pages. DOI/URL.}

\section{Awards \& Honours}
\datedline{\textbf{Academic / Research Award}, Institution}{YYYY}
\textit{Paper / Project: Optional explanatory line}\par

% Activate this section only for a genuinely developed research agenda.
\section{Current Research Agenda}
\textbf{Current Project Title}\par\vspace{1.5pt}
\puzzleline{State the genuine research puzzle or question.}
\begin{itemize}
    \item Define the substantive or empirical scope and cases.
    \item State the conceptual/theoretical distinction or mechanism being examined.
    \item State the research design, methods, and evidence strategy.
\end{itemize}

\section{Academic Leadership}
\roleentry{Academic Leadership Role}{Organisation / Institution}{YYYY--Present}
\begin{itemize}
  \item Describe academically substantive responsibilities such as programme design, research materials, training, assessment, or scholarly coordination.
\end{itemize}

\section{Professional Experience}
\roleentry{Relevant Role}{Organisation}{Mon--Mon YYYY}
\begin{itemize}
  \item Keep only details that provide relevant substantive, institutional, analytical, or research evidence.
\end{itemize}

\section{University Service}
\roleentry{Service Role}{Institution}{YYYY--YYYY}
\begin{itemize}
  \item Describe substantive service briefly.
\end{itemize}

\section{\texorpdfstring{Research Methods \& Technical Skills}{Research Methods and Technical Skills}}
\textbf{Qualitative Methods:} Method; method; method.\\[0.5pt]
\textbf{Quantitative Methods:} Method; method.\\[0.5pt]
\textbf{Developing Methods:} Only methods genuinely still being developed.\\[0.5pt]
\textbf{Software and Research Tools:} Tool; tool; tool.\\[0.5pt]
\textbf{Languages:} Language --- Level; Language --- Level.

\end{document}
